\section{Vorgehen}
Da wir beide die Fähigkeiten und das Fachwissen auf dem gleiche  Niveau besitzen, haben wir die Aufgaben und Tätigkeiten zusammen geplant und iterativ durchgeführt.
\newline

Neue Erkenntnisse haben wir uns gegenseitig beigebracht und erklärt. Bei Problemen oder Schwierigkeiten haben wir uns gegenseitig unterstützt und Lösungen gesucht.


\subsection{Code-Qualität}

Die Autoren versuchen nach dem DRY-Prinzip (Don't Repeat Yourself) zu arbeiten. Das bedeutet, dass wir versuchen, Code-Duplikate zu vermeiden und stattdessen wiederverwendbare Funktionen zu erstellen.

Der Source Code wurde auf GitHub gehostet und mit GitHub Actions automatisiert getestet. Dies ermöglicht uns, den Code in einem stabilen Zustand zu halten und die Qualität des Codes zu gewährleisten.

Anschliessend wird der Code von der anderen Person manuell geprüft. 
\newline


\subsection{Variantenentscheide}

\subsubsection{Programmiersprache: Go}

Die Autoren entschieden sich dafür, Go als Programmiersprache zu verwenden, um die Grundlagen und verschiedene Konzepte wie Go-Routines zu verstehen und anzuwenden. Zusätzlich war die Idee, Go in zukünftigen Projekten zu verwenden, sollte Go unsere Erwartungen entsprechen.

\subsubsection{Dokumentation: LaTeX}

Für LaTeX haben wir uns entschieden, da es eine verbreitete Method ist, um Dokumente zu erstellen. Im Gegensatz zu Markdown, bietet LaTeX eine Vielzahl von Funktionen, die es ermöglichen, Dokumente in einem wissenschaftlichen Format zu erstellen. \newline
Dies soll nicht Markdown ersetzen, sondern eine Alternative bieten. \newline
